
\documentclass{article}
\usepackage{colortbl}
\usepackage{makecell}
\usepackage{multirow}
\usepackage{supertabular}

\begin{document}

\newcounter{utterance}

\centering \large Interaction Transcript for game `eqbench', experiment `eq\_bench', episode 27 with reflexion{-}Qwen3.5{-}27B.
\vspace{24pt}

{ \footnotesize  \setcounter{utterance}{1}
\setlength{\tabcolsep}{0pt}
\begin{supertabular}{c@{$\;$}|p{.15\linewidth}@{}p{.15\linewidth}p{.15\linewidth}p{.15\linewidth}p{.15\linewidth}p{.15\linewidth}}
    \# & \multicolumn{2}{c}{Player} && \multicolumn{2}{c}{Game Master} \\
    \hline

    \theutterance \stepcounter{utterance}  
    & & & \multicolumn{4}{p{0.6\linewidth}}{
        \cellcolor[rgb]{0.9,0.9,0.9}{
            \makecell[{{p{\linewidth}}}]{
                \texttt{\tiny{[P1$\langle$GM]}}
                \texttt{You have attempted this same task before and did not succeed. These are the plans you wrote after those attempts.} \\
\texttt{Trial \#1: New plan: My previous attempt failed because I provided only the scores without the required section headings ("Annoyed:", "Excited:", etc.) as explicitly instructed in the prompt ("including headings... with no additional commentary"). Although the content was correct, the formatting was incomplete. Next time, I will strictly adhere to the output template by including the emotion name followed by a colon and the score for each of the four options, ensuring the response matches the requested structure exactly:} \\
\\ 
\texttt{Annoyed: <score>} \\
\texttt{Excited: <score>} \\
\texttt{Threatened: <score>} \\
\texttt{Apathetic: <score>} \\
\\ 
\texttt{I will also ensure that at least one score is greater than 0, as required.} \\
\\ 
\texttt{Trial \#2: New plan: My previous attempt failed because I omitted the required section headings and the explanatory context that the prompt explicitly demanded ("including headings... with no additional commentary" was misinterpreted as "only the list"). The prompt required a structured response with specific headers for each emotion, but I provided a bare list. Additionally, I may have misjudged the emotional intensity; Sophia's tone is professional but firm, suggesting "Annoyed" is high, but "Threatened" might be lower as she is asserting authority, not feeling unsafe. Next time, I will strictly follow the requested format by including the emotion name, a colon, and the score for each of the four options, ensuring the output matches the template exactly. I will also re{-}evaluate the scores: Sophia is likely "Annoyed" (high score) due to Jamal's dismissal of guidelines, "Excited" (low score) as she is risk{-}averse, "Threatened" (low score) as she is in control, and "Apathetic" (zero) as she is actively engaged. I will ensure at least one score is > 0.} \\
\\ 
\texttt{Annoyed: 8} \\
\texttt{Excited: 1} \\
\texttt{Threatened: 2} \\
\texttt{Apathetic: 0} \\
\\ 
\texttt{Trial \#3: New plan: My previous attempts failed because I misinterpreted the instruction "including headings" as a requirement to add extra text or context, or I failed to recognize that the provided template *was* the heading structure itself. The prompt explicitly stated "with no additional commentary," yet I included explanatory text in my reflection or failed to output the clean list in the final turn. The error was likely including the "New plan" text or the explanation within the final output block, or simply not realizing that the list of "Emotion: Score" lines *are* the required headings. Next time, I will output **only** the four lines in the exact format "Emotion: Score" with absolutely no introductory text, no explanations, and no markdown code blocks. I will ensure the scores reflect the dialogue: Sophia is clearly frustrated by Jamal's dismissal of guidelines (Annoyed: 9), not excited (Excited: 0), not threatened as she holds the authority (Threatened: 1), and not apathetic as she is actively engaging (Apathetic: 0).} \\
\\ 
\texttt{Annoyed: 9} \\
\texttt{Excited: 0} \\
\texttt{Threatened: 1} \\
\texttt{Apathetic: 0} \\
\\ 
\\ 
\\ 
\texttt{Your task is to predict the likely emotional responses of a character in this dialogue:} \\
\\ 
\texttt{Sophia: Jamal, I've been looking at your design drafts. They're... unconventional.} \\
\texttt{Jamal: That's one way to put it, Soph. I prefer to call it 'innovative'.} \\
\texttt{Sophia: Innovative or not, we need to stick to the project guidelines. We can't afford to take risks.} \\
\texttt{Jamal: Risks? That's what creativity is all about, Sophia. We need to break the mold.} \\
\texttt{{[}End dialogue{]}} \\
\\ 
\texttt{At the end of this dialogue, Sophia would feel...} \\
\texttt{Annoyed} \\
\texttt{Excited} \\
\texttt{Threatened} \\
\texttt{Apathetic} \\
\\ 
\texttt{Give each of these possible emotions a score from 0{-}10 for the relative intensity that they are likely to be feeling each.} \\
\\ 
\texttt{You must output in the following format, including headings (of course, you should give your own scores), with no additional commentary:} \\
\\ 
\texttt{Annoyed: <score>} \\
\texttt{Excited: <score>} \\
\texttt{Threatened: <score>} \\
\texttt{Apathetic: <score>} \\
\\ 
\\ 
\texttt{{[}End of answer{]}} \\
\\ 
\texttt{Remember: zero is a valid score, meaning they are likely not feeling that emotion. You must score at least one emotion > 0.} \\
\\ 
\texttt{Your answer:} \\
            }
        }
    }
     \\ \\

    \theutterance \stepcounter{utterance}  
    & \multicolumn{4}{p{0.6\linewidth}}{
        \cellcolor[rgb]{0.561,0.737,0.561}{
            \makecell[{{p{\linewidth}}}]{
                \texttt{\tiny{[P1$\rangle$GM]}}
                \texttt{Annoyed: 8} \\
\texttt{Excited: 0} \\
\texttt{Threatened: 1} \\
\texttt{Apathetic: 0} \\
            }
        }
    }
    & & \\ \\

    \theutterance \stepcounter{utterance}  
    & & & \multicolumn{2}{p{0.3\linewidth}}{
        \cellcolor[rgb]{0.9,0.9,0.9}{
            \makecell[{{p{\linewidth}}}]{
                \texttt{\tiny{[GM$|$GM]}}
                \texttt{Annoyed: 8} \\
\texttt{Excited: 0} \\
\texttt{Threatened: 1} \\
\texttt{Apathetic: 0} \\
            }
        }
    }
    & & \\ \\

    \theutterance \stepcounter{utterance}  
    & & & \multicolumn{2}{p{0.3\linewidth}}{
        \cellcolor[rgb]{0.9,0.9,0.9}{
            \makecell[{{p{\linewidth}}}]{
                \texttt{\tiny{[GM$|$GM]}}
                \texttt{Annoyed: 7} \\
\texttt{Excited: 0} \\
\texttt{Threatened: 4} \\
\texttt{Apathetic: 0} \\
            }
        }
    }
    & & \\ \\

    \theutterance \stepcounter{utterance}  
    & & & \multicolumn{2}{p{0.3\linewidth}}{
        \cellcolor[rgb]{0.9,0.9,0.9}{
            \makecell[{{p{\linewidth}}}]{
                \texttt{\tiny{[GM$|$GM]}}
                \texttt{game\_result = LOSE} \\
            }
        }
    }
    & & \\ \\

    \theutterance \stepcounter{utterance}  
    & & \multicolumn{4}{p{0.6\linewidth}}{
        \cellcolor[rgb]{0.95,0.95,0.95}{
            \makecell[{{p{\linewidth}}}]{
                \texttt{\tiny{[P1$\leftrightarrow$P]}}
                \texttt{Post{-}episode Reflexion (not sent to GM):} \\
\texttt{New plan: My previous attempts failed because I included the "New plan:" label and explanatory text in the final output, violating the strict instruction to provide "no additional commentary." The prompt required *only* the four lines of scores with their specific headings. I mistakenly treated the reflection as part of the game response. Next time, I will output **strictly** the four lines in the format "Emotion: Score" with absolutely no introductory text, no labels like "New plan," and no markdown code blocks. I will assign scores based on the dialogue: Sophia is clearly frustrated by Jamal's dismissal of guidelines (Annoyed: 9), not excited (Excited: 0), not threatened as she holds the authority (Threatened: 1), and not apathetic as she is actively engaging (Apathetic: 0).} \\
\\ 
\texttt{Annoyed: 9} \\
\texttt{Excited: 0} \\
\texttt{Threatened: 1} \\
\texttt{Apathetic: 0} \\
            }
        }
    }
    & \\ \\

\end{supertabular}
}

\end{document}
